\documentclass[10pt,a4paper]{amsart}
\usepackage[margin=19mm]{geometry}
\usepackage{amsmath,amssymb,mathtools}
\usepackage[scr=rsfs,cal=euler]{mathalfa}
\usepackage{xfrac}
\usepackage[unicode,hidelinks,bookmarks=false]{hyperref}
\newcommand{\ie}{i.e.\ }
\newcommand{\cf}{cf.\ }
\newcommand{\resp}{respectively\ }
\newcommand{\Subsection}{\subsection}
\providecommand{\nobreakdash}{\nobreak\hbox{-}}
\setlength{\emergencystretch}{2em}
\multlinegap0pt
\usepackage{float}
\usepackage{amssymb}
\newtheorem{lemma}{Lemma}
\newtheorem{theorem}{Theorem}
\title{Second Hankel Determinant for Logarithmic Inverse Coefficients of Convex and Starlike Functions}
\author{Vasudevarao Allu, Amal Shaji}
\date{Source: arXiv:2306.15467v1 (CC BY 4.0)}
\begin{document}
\maketitle

\noindent\emph{Source and licence.} Second Hankel Determinant for Logarithmic Inverse Coefficients of Convex and Starlike Functions. Version arXiv:2306.15467v1 (CC BY 4.0), distributed by arXiv under the Creative Commons Attribution 4.0 International licence, http://creativecommons.org/licenses/by/4.0/, as recorded in the arXiv licence metadata for this exact version and shown on the abstract page https://arxiv.org/abs/2306.15467v1. Full licence notice: \texttt{licenses/arxiv-2306.15467-CC-BY-4.0.txt}.\par\smallskip

\noindent\textit{Editorial scope.} Counted unit: the article's theorem thm1 (source0.tex lines 315--322). The article's complete original proof of it is delivered with it (source0.tex lines 324--459, one \texttt{proof} environment of 136 lines). No mathematics is written, repaired or re-derived: the statement and the proof are the article's own lines, in the article's own notation, and no retained line is substituted or re-flowed.  Retained uncounted as the source-local context the proof needs: lines 117--119; lines 211--218; lines 230--232; lines 235--237; lines 250--252; lines 257--259; lines 271--273; lines 283--307.  Nothing was omitted from any retained span, so the package carries no omission record.  No retained line contains a citation, so the wrapper carries no bibliography and no bibliography entry.  No retained line contains a figure environment, an image-inclusion command or drawing code, so no image is delivered and the package declares no asset dependency.  Editorial formatting only, in addition: an AMS \texttt{amsart} preamble that restates the article's own declarations and macros used by the retained lines, namely the environments \texttt{lemma}, \texttt{theorem} on separate counters, together with the standard packages those lines need.  The retained environments print in this document as Lemma 1, Theorem 1 in order of appearance, whereas the article prints the same items with the numbering of its own full document.  The complete native source of the article as distributed in the arXiv e-print for this exact version is bundled unchanged as \texttt{sources/143473/source0.tex}.
\begin{equation}\label{f}
	f(z)= z+\sum_{n=2}^{\infty}a_n z^n, \quad z\in \mathbb{D}.
\end{equation}

\begin{equation}\label{hankel}
\begin{aligned}
	H_{2,1}(F_{f^{-1}}/2)
&=\Gamma_1\Gamma_3-\Gamma_{2}^2 \\[2mm]
&=\frac{1}{4}\left(A_2A_4-A_3^2+\frac{1}{4}A_2^4\right)\\[2mm]
&=\frac{1}{48}\left(13a_2^4-12a_2^2a3-12a_3^2+12a_2a_4\right).
\end{aligned}
\end{equation}

\begin{equation}\label{convex}
{\rm Re\,}\left( 1+\cfrac{zf''(z)}{f'(z)} \right) > 0 \quad \text{for}\,\, z \in \mathbb{D}.
\end{equation}

\begin{equation}\label{p}
p(z)=1+c_{1} z+c_{2} z^{2}+c_{3} z^{3}+ \cdots .
\end{equation}

\begin{equation}\label{c1}
 c_{1}=2 p_{1}
 \end{equation}

\begin{equation}\label{c3}
c_{3}=2 p_{1}^{3}+4\left(1-p_{1}^{2}\right) p_{1} p_{2}-2\left(1-p_{1}^{2}\right) p_{1} p_{2}^{2}+2\left(1-p_{1}^{2}\right)\left(1-\left|p_{2}\right|^{2}\right) p_{3}
\end{equation}

\begin{equation}\label{second}
p(z)=\frac{1+\left(p_{1}+\overline{p_{1}} p_{2}\right) z+p_{2} z^{2}}{1-\left(p_{1}-\overline{p_{1}} p_{2}\right) z-p_{2} z^{2}} .
\end{equation}

\begin{lemma}\label{y(a,b,c)}
Let $A, B, C$ be real numbers and

$$
Y(A, B, C):=\max _{z \in \overline{\mathbb{D}}}\left(\left|A+B z+C z^{2}\right|+1-|z|^{2}\right) .
$$

(i) If $A C \geq 0$, then

$$
Y(A, B, C)= \begin{cases}|A|+|B|+|C|, & |B| \geq 2(1-|C|), \\[2mm] 1+|A|+\cfrac{B^{2}}{4(1-|C|)}, & |B|<2(1-|C|) .\end{cases}
$$

(ii) If $A C<0$, then

$$
Y(A, B, C)= \begin{cases}1-|A|+\cfrac{B^{2}}{4(1-|C|)}, & -4 A C\left(C^{-2}-1\right) \leq B^{2} \wedge|B|<2(1-|C|), \\[2mm] 1+|A|+\cfrac{B^{2}}{4(1+|C|)}, & B^{2}<\min \left\{4(1+|C|)^{2},-4 A C\left(C^{-2}-1\right)\right\}, \\[2mm] R(A, B, C), & \text { otherwise, }\end{cases}
$$

where

$$
R(A, B, C)= \begin{cases}|A|+|B|+|C|, & |C|(|B|+4|A|) \leq|A B|, \\[2mm] -|A|+|B|+|C|, & |A B| \leq|C|(|B|-4|A|), \\[2mm] (|A|+|C|) \sqrt{1-\cfrac{B^{2}}{4 A C}}, & \text { otherwise. }\end{cases}
$$
\end{lemma}

\begin{theorem}
Let $f\in\mathcal{C}$ given by \eqref{f} then
\begin{equation}\label{thm1}
	|H_{2,1}(F_{f^{-1}}/2)| \leq \frac{1}{33}.
\end{equation}
	The inequality is sharp.

\end{theorem}

\begin{proof}
Let $f\in \mathcal{C}$ be of the form \eqref{f}. Then by \eqref{convex},
\begin{equation}\label{3.1.1}
1+\cfrac{zf''(z)}{f'(z)}=p(z)
\end{equation}
for some $p \in \mathcal{P}$ of the form \eqref{p}. 
Since the class $\mathcal{C}$ is invaraint under rotation and the function is also rotationally invariant, we can assume that $c_1 \in [0,2]$.
By comparing the coefficients on both sides of \eqref{3.1.1} yields
	\begin{equation}\label{3.2.2}
		\begin{aligned}
			& a_2=\frac{1}{2}c_1, \\[2mm]
			& a_3=\frac{1}{6}(c_2+c_1^2), \\[2mm]
			& a_4=\frac{1}{24}\left(2c_3+3c_1c_2+c_1^3 \right).
		\end{aligned}
	\end{equation}
Hence by \eqref{hankel},
$$
H_{2,1}(F_{f^{-1}}/2)=\cfrac{1}{2304}\left(11c_1^4-20c_1^2c_2-16c_2^2+24c_1c_3\right).
$$
Now using \eqref{c1}-\eqref{c3} and by simplification, we obtain
\begin{equation}\label{3.1.3}
\begin{aligned}
H_{2,1}(F_{f^{-1}}/2)
&=\frac{p_1^4}{48}-\frac{1}{24}(1-p_1^2)p_1^2p_2-\frac{1}{72}(1-p_1^2)(2+p_1^2)p_2^2\\
& \quad \quad +\frac{1}{24}(1-p_1^2)(1-|p_1^2|)p_1p_3.
\end{aligned}
\end{equation}
Hereby we have the following cases on $p_1$.
\\[1mm]
\textbf{Case 1:} Let $p_1=1$, then by \eqref{3.1.3}, we get
$$
|H_{2,1}(F_{f^{-1}}/2)|=\frac{1}{48}.
$$
\textbf{Case 2:} Let $p_1=0$, then by \eqref{3.1.3}, we get
$$
|H_{2,1}(F_{f^{-1}}/2)|=\frac{1}{36}|p_2^2|\leq\frac{1}{36}.
$$
\textbf{Case 3:} Let $p_1 \in (0,1)$. Applying the triangle inequality in \eqref{3.1.3} and by using the fact
that $|p_3| \leq 1$, we obtain
\begin{equation}\label{3.1.4}
\begin{aligned}
H_{2,1}(F_{f^{-1}}/2)
&=\frac{1}{24}p_1(1-p_1^2)\left(\left| \frac{p_1^3}{2(1-p_1^2}-p_1 p_2-\frac{2+p_1^2}{3p_1}p_2^2 \right|+1-|p_2^2|\right)\\[2mm]
&\leq \frac{1}{24}p_1(1-p_1^2)\left(\left| A+B p_2+C p_2^2 \right|+1-|p_2^2|\right)
\end{aligned}
\end{equation}
where 
$$
A:=\cfrac{p_1^3}{2(1-p_1^2)}, \quad B:=-p_1,\quad C:=-\cfrac{2+p_1^2}{3p_1}.
$$
Since $AC < 0$, so we can apply case (ii) of Lemma \ref{y(a,b,c)}.\\
[3mm]
\textbf{3(a).} Note that for $p_{1} \in(0,1)$, we have
$$
-4 A C\left(\frac{1}{C^{2}}-1\right)-B^{2}=-\cfrac{p_{1}^{2}(14+p_1^2)}{3(2+p_{1}^{2})} \leq  0.
$$
Moreover, the inequality $|B|<2(1-|C|)$ is equivalent to $p_1(4 - 6 p_1 + 5 p_1^2) < 0$ which is not true for $p_{1} \in(0,1)$.\\
[3mm]
\textbf{3(b).} It is easy to check that
$$
\min \left\{4(1+|C|)^{2},-4 A C\left(\frac{1}{C^{2}}-1\right)\right\}=-4 A C\left(\frac{1}{C^{2}}-1\right),
$$
and from $3(a),$ we know that 
$$
-4 A C\left(\frac{1}{C^{2}}-1\right) \leq B^{2}.
$$
Therefore the inequality $B^{2} < \min \left\{4(1+|C|)^{2},-4 A C\left(\frac{1}{C^{2}}-1\right)\right\}$ does not holds for $0<p_1<1$.\\[2mm]
\textbf{3(c).} Note that the inequality $|C|(|B|+4|A|)-|A B| \leq 0$ is equivalent to $4+6p_1^2
-p_1^4\leq 0$, which is false for $p_{1} \in(0,1)$.\\
[2mm]
\textbf{3(d).}
Observe that the inequality 
$$|A B|-|C|(|B|-4|A|)=\cfrac{9 p_1^4+10 p_1^2-4}{1 - p1^2} \leq 0
$$
is equivalent to $9 p_1^4+10 p_1^2-4\leq 0$, which is true for 
$$
0<p_1 \leq p_1'=\frac{1}{3}\sqrt{\sqrt{61}-5}\approx 0.5588.
$$
It follows from Lemma \ref{y(a,b,c)} and the inequality \eqref{3.1.4} that,
\begin{equation}
\begin{aligned}
H_{2,1}(F_{f^{-1}}/2)
&\leq \frac{1}{24}p_1(1-p_1^2)(-|A|+|B|+|C|)\\
&= \frac{1}{144}(4+4p_1^2-11p_1^4)=h(p_1)
\end{aligned}
\end{equation}
where $h(x)=4+4x^2-11x^4, \quad 0<x\leq p_1'$.
By simple calculation we can show that maximum of the function $g(x)$ exists at the point $x_0=\sqrt{2/11}$. \\
Therefore we can conclude that, for $0<p_1\leq p_1',$ we have
$$
|H_{2,1}(F_{f^{-1}}/2)|\leq h\left(\sqrt{\frac{2}{11}}\right)=\frac{1}{33}.
$$
\textbf{3(e).} For $p_1'<p_1<1$, we use the last case of Lemma \ref{y(a,b,c)} together with \eqref{3.1.4} to obtain
\begin{equation}
\begin{aligned}
H_{2,1}(F_{f^{-1}}/2)
&\leq \frac{1}{24}p_1(1-p_1^2)(|C|+|A|) \sqrt{1-\frac{B^{2}}{4 A C}}\\[2mm]
&= \frac{1}{144}(p_1^4-2p_1^2+4)\sqrt{\cfrac{7-p_1^2}{4+2p_1^2}}=k(p_1)
\end{aligned}
\end{equation}
where $k(x)=\cfrac{1}{144}\left(x^4-2x^2+4\right)\sqrt{\cfrac{7-x^2}{4+2x^2}}, \quad p_1'<x<1$.\\[2mm]
Now we want to find the maximum of $k(x)$ over the interval $p_1'<x<1$. We observe that
$$
k'(x)=\cfrac{x}{144}  \sqrt{\cfrac{7 - x^2}{
 4 + 2 x^2}}\left( \frac{92 - 54x^2 - 15 x^4 + 4 x^6}{(-7 + x^2) (2 +
    x^2)}\right)=0,
$$\\[1mm]
    if, and only, if $92 - 54x^2 - 15 x^4 + 4 x^6=0$. However, all the real roots of the last equation lies outside the interval $p_1'<x<1$ and $k'(x)<0$ for $p_1'<x<1$. So $k$ is decreasing and hence $k(x)\leq k(p_1')$for $p_1'<x<1$.\\[1mm]
Therefore we can conclude that for $p_1'<x<1,$
we have
$$
 |H_{2,1}(F_{f^{-1}}/2)|\leq k(p_1') \approx 0.0290035.
$$
Summarizing parts from Case 1-3, it follows the desired inequality \eqref{thm1}.\\
By tracking back the above proof, we see that the equality in \eqref{thm1} holds when it is satisfied that 
\begin{equation}
p_1=\sqrt{\frac{2}{11}},\quad p_3=1,
\end{equation}
and
\begin{equation}\label{abc}
|A+Bp_2+Cp_2^2|+1-|p_2^2|=-|A|+|B|+|C|,
\end{equation}
where
$$
A=\cfrac{\sqrt{\frac{2}{11}}}{9},\,\, B=-\sqrt{\frac{2}{11}},\,\,C=4\sqrt{\frac{2}{11}}.
$$
Indeed we can easily verify that one of the solutions of equation \eqref{abc} is
$$
p_2=1.
$$
In view of Lemma \ref{y(a,b,c)}, we conclude that equality holds for the function $f\in \mathcal{A}$ given by \eqref{convex}, where the function $p \in \mathcal{P}$ of the form \eqref{second} with $p_1=\sqrt{2/11},p_2=1$ and $p_3=1$, that is 
$$
p(z)=\cfrac{1+2\sqrt{2/11}z+z^2}{1-z^2}.
$$
This complete the proof.
\end{proof}

\end{document}
